% chapter-algorithm.tex — \input'd from main.tex
%
% FOUR AUDIENCES — every line must work for all four simultaneously.
%
% AUDIENCE 1: Human readers (Jörn, Kai, Elizabeth).
%   - Jörn: knows enough symplectic geometry, optimization, this entire thesis, the bibliography materials
%   - Kai: knows a lot of symplectic geometry
%   - Elizabeth: knows a lot of optimization
%   - All three: know typical basics such as linear algebra, polytopes, analysis, ...
%   - Want the main result (algorithm + theorem) upfront on page 1, so they can
%     read it and stop if that's all they need.
%   - Will skim definitions and revisit if confused.
%   - All proofs are skippable: clearly marked start/end, nothing outside a proof
%     references proof internals.
%   - Value: the algorithm itself, the interesting proof ideas, geometric intuition.
%   - Remarks, examples, and illustrations live next to relevant lemmas/proofs.
%     Redundancy is fine — readers pick what they want.
%
% AUDIENCE 2: Imaginary master student (nominal target audience of a master thesis).
%   - Knows linear algebra, analysis, basic topology, intro symplectic geometry,
%     intro optimization — a typical math master student's background.
%   - Every definition must be stated in full, not deferred to the literature.
%   - Must be able to follow the chapter linearly without external references.
%   - This audience is why there's motivation even for standard definitions, while
%     the other three audiences are familiar with the literature already.
%
% AUDIENCE 3: QC agents (pedantic correctness verification, AI).
%   - Verify one chunk at a time, trusting previously verified chunks.
%   - Each chunk is self-contained: a QC agent can read [definitions]
%     [propositions] [lemma X] [theorem Y] [proof of Y] and check [proof of Y]
%     while trusting that the rest is already verified.
%   - For every proof step: if the agent cannot immediately confirm "yes, that
%     follows directly," the step fails review.
%   - If any step is incorrect, the chunk fails review.
%   - Words must have clear, specific meanings. Natural language errs towards
%     explicitness to prevent misunderstandings or overlooked edge cases.
%   - We never state anything incorrect, even if the incorrectness is not fatal.
%
% AUDIENCE 4: Downstream agents (Rust implementers, test writers, AI).
%   - Need full detail in all definitions, lemmas, proofs, algorithm specs.
%   - Need all properties listed — including properties not cited later in
%     the chapter — because they generate unit tests and proptests from them.
%     E.g., capacity axioms: useful for intuition AND for test cases.
%   - Need concrete values and example calculations that illustrate the
%     algorithms and proof steps.
%
% CONTENT RULES:
%   1. Every definition, theorem statement, and proof strategy must be traceable
%      to a specific instruction from Jörn in the chat.
%      Do NOT write anything Jörn did not dictate.
%   2. No definition or theorem statement may be deferred to the literature.
%      Every definition is stated in full. Every theorem is stated in full.
%   3. A proof MAY be deferred to the literature ONLY IF:
%      (a) the theorem exceeds thesis scope due to complexity, AND
%      (b) the proof is not relevant to the thesis — only the theorem is.
%   4. Notation and definitions must match correspondence.tex exactly.
%      If correspondence.tex uses symbol X, this file uses symbol X.
%      No synonyms, no alternative forms, no "equivalent" restatements
%      unless a lemma proving equivalence is included.
%
% FORMAT RULES:
%   - \definition, \lemma, \theorem, \proposition, \proof environments for
%     mathematical content. \remark and \example for context/intuition/illustrations.
%   - No prose paragraphs outside environments (except minimal connective text).
%   - All calculations displayed as formulas, never described in English sentences.
%   - Emphasis proportional to importance. No full sections for trivial identities.
%
% QC PROCEDURE:
%   A QC agent is given correspondence.tex and this file. For each environment:
%   (a) Does the definition/statement match correspondence.tex? Find one deviation.
%   (b) Does each proof step follow immediately from stated assumptions? Find one gap.
%   If the QC agent cannot find a deviation or gap, the environment passes.

%%% =================================================================
%%% Section 1: Main Result
%%% =================================================================

\section{Main Result}\label{sec:main-result}

% QC: polytope per Definition~\ref{def:polytope}, facet data per Definition~\ref{def:facets},
% capacity per Definition~\ref{def:ehz-capacity}.

\begin{theorem}[EHZ capacity formula {\cite{hk2017}}]\label{thm:main}
Let $K \subset \mathbb{R}^4$ be a polytope
with $F$ facets, outward unit normals $n_1, \ldots, n_F \in S^3$,
and heights $h_1, \ldots, h_F > 0$.
Then
\begin{equation}\label{eq:ehz-main}
  c_{\mathrm{EHZ}}(K) = \frac{1}{2}
  \left[
    \max_{\substack{
      \sigma \in S_F \\
      \beta \in M(K)
    }}
    \sum_{1 \leq j < i \leq F}
      \beta_{\sigma(i)} \beta_{\sigma(j)} \,
      \omega_0(n_{\sigma(i)}, n_{\sigma(j)})
  \right]^{-1},
\end{equation}
where
\begin{equation}\label{eq:constraint-set}
  M(K) = \left\{
    (\beta_i)_{i=1}^{F} :
    \beta_i \geq 0, \quad  % not >: M(K) includes boundary where some \beta_i = 0
    \sum_{i=1}^{F} \beta_i h_i = 1, \quad
    \sum_{i=1}^{F} \beta_i n_i = 0
  \right\}.
\end{equation}
\end{theorem}

% Proof: Section~\ref{sec:simple-minimizer} (search space)
% and Section~\ref{sec:algorithm-correctness} (algorithm correctness).

\medskip\noindent
\textbf{Algorithm.}\label{alg:ehz}
The following algorithm computes~\eqref{eq:ehz-main}.

\emph{Input:} $(n_1, h_1), \ldots, (n_F, h_F)$.
\emph{Output:} $c_{\mathrm{EHZ}}(K)$.

\medskip\noindent
For each nonempty subset $S \subseteq \{1, \ldots, F\}$
and each permutation $\sigma$ of $S$
(identified up to cyclic shifts):

\begin{enumerate}
\item \textbf{Define.}
  The normal matrix $N \in \mathbb{R}^{|S| \times 4}$,
  action matrix $H \in \mathbb{R}^{|S| \times |S|}$,
  and height vector $\eta \in \mathbb{R}^{|S|}$:
  \begin{align}
    N_{id} &= n_{\sigma(i),d},
      & i &= 1, \ldots, |S|, \quad d = 1, \ldots, 4, \\
    H_{ij} &= \begin{cases}
      \omega_0(n_{\sigma(i)}, n_{\sigma(j)}) & \text{if } i < j, \\
      0 & \text{if } i = j, \\
      \omega_0(n_{\sigma(j)}, n_{\sigma(i)}) & \text{if } i > j,
    \end{cases} \\
    \eta_i &= h_{\sigma(i)},
      & i &= 1, \ldots, |S|.
  \end{align}

\item \textbf{Solve.}
  \begin{equation}\label{eq:linear-system}
    \begin{pmatrix} N^\top & 0 \\ H & -N \\ \eta^\top & 0 \end{pmatrix}
    \begin{pmatrix} \beta \\ \lambda \end{pmatrix}
    = \begin{pmatrix} 0 \\ 0 \\ 1 \end{pmatrix},
    \qquad
    \beta \in \mathbb{R}^{|S|}, \quad \lambda \in \mathbb{R}^4.
  \end{equation}

\item \textbf{Filter.}
  Discard $(S, \sigma)$ if $\beta_i \leq 0$ for any $i$.
  % not <: we also discard \beta_i = 0 (degenerate); the remark below requires \beta_i > 0

  \begin{remark}\label{rem:critical-point}
  The solution $\beta$ with $\beta_i > 0$
  is the critical point of
  $Q(\beta) = \sum_{1 \leq j < i \leq |S|}
    \beta_i \beta_j \, \omega_0(n_{\sigma(i)}, n_{\sigma(j)})$
  subject to $\sum_{i} \beta_i n_{\sigma(i)} = 0$,
  $\sum_{i} \beta_i h_{\sigma(i)} = 1$, and $\beta_i > 0$.
  \end{remark}

\item \textbf{Evaluate.}
  \begin{equation}\label{eq:action-formula}
    A(S, \sigma) = \tfrac{1}{2} \, Q(\beta)^{-1}.
  \end{equation}
\end{enumerate}

\noindent
Return $c_{\mathrm{EHZ}}(K) = \min_{S,\sigma} A(S, \sigma)$.

%%% =================================================================
%%% Section 2: Definitions
%%% =================================================================

\section{Definitions}\label{sec:definitions}

% --- Smooth setting ---

\begin{definition}[Standard symplectic form]\label{def:symplectic-form}
\end{definition}

\begin{definition}[Liouville form]\label{def:liouville-form}
\end{definition}

\begin{definition}[Convex body]\label{def:convex-body}
\end{definition}

\begin{definition}[Support function]\label{def:support-function}
\end{definition}

\begin{definition}[Gauge function]\label{def:gauge-function}
\end{definition}

\begin{definition}[Hamiltonian]\label{def:hamiltonian}
\end{definition}

\begin{definition}[Normal cone]\label{def:normal-cone}
\end{definition}

\begin{definition}[Generalized Reeb orbit --- smooth]\label{def:reeb-orbit-smooth}
\end{definition}

\begin{definition}[Action of an orbit]\label{def:action}
\end{definition}

\begin{definition}[EHZ capacity]\label{def:ehz-capacity}
\end{definition}

\begin{lemma}[Symplectic capacity axioms]\label{lem:capacity-axioms}
\end{lemma}

% --- Polytope setting ---

\begin{definition}[Polytope]\label{def:polytope}
% 0 in int(K) — Jörn to confirm where this is used
\end{definition}

\begin{definition}[Facets, normals, heights]\label{def:facets}
\end{definition}

\begin{definition}[Facet Reeb vectors]\label{def:reeb-vectors}
\end{definition}

\begin{definition}[Generalized Reeb orbit --- polytope]\label{def:reeb-orbit-polytope}
\end{definition}

%%% =================================================================
%%% Section 3: Simple Minimizer Existence
%%% =================================================================

\section{Simple Minimizer Existence}\label{sec:simple-minimizer}

\begin{theorem}[Existence of Reeb orbits, positive infimum, minimum attained]\label{thm:orbit-existence}
% Cited, no proof. Smooth: established. Polytope: CH2021?
\end{theorem}

\begin{theorem}[Simple minimizer structure]\label{thm:simple-minimizer}
% Piecewise linear, pure facet Reeb vectors, contiguous intervals
\end{theorem}

%%% =================================================================
%%% Section 4: Algorithm Correctness
%%% =================================================================

\section{Algorithm Correctness}\label{sec:algorithm-correctness}

\begin{definition}[Dual variable]\label{def:dual-variable}
\end{definition}

\begin{definition}[Dual functional]\label{def:dual-functional}
\end{definition}

\begin{definition}[$Q$-function]\label{def:q-function}
\end{definition}

\begin{definition}[Constraint set]\label{def:constraint-set}
\end{definition}

\begin{definition}[Action matrix]\label{def:action-matrix}
\end{definition}

\begin{theorem}[Optimization formulation]\label{thm:optimization}
% HK2017 max formula
\end{theorem}

\begin{lemma}[Domain decomposition]\label{lem:domain-decomposition}
% Union of open domains, parameterized by (subset of facets, cyclic order)
\end{lemma}

\begin{lemma}[Global maximum at local maximum]\label{lem:global-at-local}
\end{lemma}

\begin{corollary}[Enumeration computes capacity]\label{cor:enumeration}
\end{corollary}

\begin{lemma}[Local maximum characterization]\label{lem:local-max}
% Explicit linear equation, saddle point check, domain membership check
\end{lemma}

% Algorithm (verbatim copy from Section 1)

